\chapter{The Generative AI Layer}
\label{ch:ai}

This chapter documents the part of the system that is genuinely about generative AI:
how a turn is classified, how knowledge is retrieved and cited, how the hint ladder and
the Auditor are built to bound model usage, and how the whole layer degrades when the
model is not there. The organising principle throughout is that \emph{the model is the
most expensive and least reliable component}, so every use of it is a deliberate budgeted
decision with a template-based fallback.

\section{Group Pulse: one call, many decisions}

Per learner turn, the system needs four pieces of information: the intent of the message,
the language and register it was written in, whether it contains unsafe content or PII,
and how difficult the underlying question is. The naive implementation makes four calls.
HALQA makes one schema-constrained compact call and shares its result across four
consumers.

\begin{table}[htbp]
\centering
\caption{The Group Pulse schema and its consumers. One compact call per turn window serves
credits, confusion tracking, register detection and the difficulty gate.}
\label{tab:pulse}
\footnotesize
\begin{tabularx}{\textwidth}{@{}p{3.1cm}p{3.0cm}X@{}}
\toprule
\textbf{Field} & \textbf{Domain} & \textbf{Consumed by} \\
\midrule
\code{intent} & \code{ask}, \code{explain}, \code{help\_offer}, \code{dispute},
  \code{chitchat}, \code{meta} & Router (task selection) and credit weights
  (\S\ref{sec:credits}) \\
\code{helpful\_0\_2} & $0..2$ & Credit award for explanations; S-B-U \emph{Build} ring \\
\code{bridging\_bool} & boolean & Connects two learners' ideas; \emph{Bridge Builder} badge \\
\code{skeptic\_bool} & boolean & Kind-skeptic credit; disagreement is data, not shame \\
\code{confusion\_0\_1} & $0..1$ & Confusion signal; drives hint eagerness and
  \code{needs\_human} \\
\code{lang} & \code{ar}, \code{darija}, \code{fr}, \code{en} & Response register and
  retrieval language filter \\
\bottomrule
\end{tabularx}
\end{table}

Two implementation details matter for correctness under latency. First, the response is
parsed through a narrowing function that validates every field and substitutes a safe
default (\code{chitchat}, neutral helpfulness) rather than throwing: a malformed model
reply must degrade one turn, not fail the request. Second, the call is guarded so that
\emph{at most one} Group Pulse runs per turn window; a burst of messages collapses into a
single classification. Both are unit-tested in \code{router-pure.spec.ts}.

\section{Retrieval-augmented generation}
\label{sec:rag}

\subsection{The knowledge base}

Knowledge lives in versioned JSON seed files rather than in a vector store, because the
corpus is small, curated and needs to be auditable. This submission ships two packs ---
\emph{binary search} and \emph{fractions} --- totalling 39 chunks, each with a stable
identifier, a language, a kind, and an explicit list of \emph{claims} it supports.

\begin{table}[htbp]
\centering
\caption{Knowledge base composition. Every chunk carries claim checklists, which is what
lets the Auditor verify a peer answer without a second retrieval pass.}
\label{tab:kb}
\footnotesize
\begin{tabular}{@{}lrrrr@{}}
\toprule
\textbf{Pack} & \textbf{Chunks} & \textbf{ar} & \textbf{darija} & \textbf{fr / en} \\
\midrule
\code{binary\_search.json} & 27 & 6 & 6 & 15 \\
\code{fractions.json}      & 12 & 1 & 1 & 10 \\
\midrule
\textbf{Total}             & \textbf{39} & 7 & 7 & 25 \\
\bottomrule
\end{tabular}
\end{table}

The manifest \code{kb/SOURCES.md} lists every chunk identifier with its provenance and
licence, and a test asserts that the seed files and the manifest agree in both
directions --- a chunk listed but not shipped, or shipped but not listed, fails the build.
All content is team-written and synthetic, which also means the data sent to the provider
contains no third-party material and no learner data.

\subsection{Index and query}

Chunks are indexed at boot with TF-IDF over a normalised token field. The index builds in
under a second, which is what allows it to be rebuilt on every process start with no
persistence layer and no warm-up window.

Each grounded query is rewritten into three template variants (base, base with the level
appended, base with an example cue appended), every variant is scored, results are
deduplicated by chunk identifier keeping the best score, ties are broken deterministically
by identifier, and the top five are attached to the prompt together with their claim
checklists.

\begin{lstlisting}[language=Java,caption={Query across the three rewrite variants, with deterministic dedupe.}]
search(intent, level, sessionLang, topK = 5) {
  const norm = normaliseArabizi(intent);
  const best = new Map();
  for (const variant of rewriteQuery(norm, level)) {
    for (const hit of this.index.search(variant, sessionLang, topK)) {
      // keep the best score per chunk, so expansion can only help recall
      const prev = best.get(hit.chunkId);
      if (!prev || hit.score > prev.score) best.set(hit.chunkId, hit);
    }
  }
  return [...best.values()]
    .sort((a, b) => b.score - a.score || (a.chunkId < b.chunkId ? -1 : 1))
    .slice(0, topK);
}
\end{lstlisting}

Determinism here is not aesthetic. The demo, the evals and the smoke test all compare
outputs across runs, so retrieval that varied between runs would make every downstream
assertion flaky. Sorting ties by identifier also means the escalation gate --- which counts
\emph{conflicting} chunks --- sees a stable set.

\subsection{Mandatory citations}

Every grounded answer must cite chunk identifiers, and the interface renders them as
citation chips beneath the answer. This is a stronger requirement than it looks: a
citation is only useful if the claim it supports is checkable, so the same claim
checklists that are attached to the prompt are stored with the answer and later used by
the Auditor to verify peer claims. One artefact, three consumers --- grounding,
verification and the offline pack.

\begin{note}[Why citations reduce load rather than increase it]
An answer with five citations is longer, but it is also more likely to be accepted without
a follow-up question, which is the expensive path. In practice the citation obligation is a
token cost paid once to avoid a second turn, and it makes the Auditor's job a set
comparison rather than a fresh generation.
\end{note}

\subsection{Arabizi normalisation}
\label{sec:arabizi}

Tunisian learners routinely write Darija in Latin script with digits standing in for
letters: \code{3} for \emph{ain}, \code{7} for \emph{ha}, \code{9} for \emph{qaf}, and
\code{5} or \code{kh} for \emph{kha}. Two learners can ask the same question in Arabic
script and in Arabizi, and a tokenizer that treats them as unrelated strings will fail to
retrieve the same chunk for both.

The domain layer therefore ships a pure normaliser and a register detector. Detection runs
before retrieval so that a chunk in the wrong register is never attached to a prompt: a
Darija question receives Darija chunks. Because the functions are pure and
dependency-free, they are tested on a fixed corpus of real orthographic variants, and the
same code path is reused by the cache key builder so that an equivalent re-ask in a
different encoding can still hit the cache.

\section{The hint ladder}
\label{sec:hints}

The ladder is the mechanism by which the system \emph{scaffolds rather than solves}. Four
levels escalate from a nudge to a full worked solution, and the passport decides how early
the ladder may begin for a given learner.

\begin{table}[htbp]
\centering
\caption{The four hint levels, all served at zero model cost except the last.}
\label{tab:ladder}
\footnotesize
\begin{tabularx}{\textwidth}{@{}p{1.0cm}p{2.5cm}X p{3.3cm}@{}}
\toprule
\textbf{Level} & \textbf{Kind} & \textbf{Content} & \textbf{Source} \\
\midrule
L1 & Nudge & A question that names the missing precondition without stating it:
   \emph{``What must be true about the array before you start halving?''} & T0 template \\
L2 & Analogy & A concrete counter-example that breaks the learner's current mental model:
   \emph{``Is the list [5,1,4] sorted? Try one halving step and see what breaks.''}
   & T0 template \\
L3 & Worked start & The first step of a \emph{different} instance, so the learner
   transfers the move rather than copying the answer & T0 template \\
L4 & Full solution & The hidden Quick Take, with citations & Prepared during the peer
   window \\
\bottomrule
\end{tabularx}
\end{table}

L1--L3 are templates keyed by topic and language. They cost nothing, they cannot fail, and
they are available in every degradation state including \code{BLACK}. L4 is the prepared
answer, and it is \emph{leak-gated in code}: the service refuses L4 unless the learner has
already advanced the hint state to at least level three or pressed skip. The guard is a
condition on persisted state, not an instruction in a prompt, and a unit test asserts the
refusal:

\begin{lstlisting}[language=Java,caption={The L4 leak gate. A prompt-only instruction would be advisory; this is not.}]
// T0 leak gate (enforced in code, not prompt): L4 requires hintState >= 3
if (level === 4 && !(q.hintState >= 3 || q.skipPressed)) throw Err.denied('HINT_L4_LOCKED');
\end{lstlisting}

The gate exists because a ladder that can be short-circuited is not a ladder. If a learner
could request L4 immediately, the scaffolding would be decorative and the system would
revert to an answer machine.

\section{The Auditor}
\label{sec:auditor}

Peers first creates a new failure mode: the group can now converge confidently on a wrong
answer. The Auditor exists to prevent that, and its design is shaped by the quota
constraint --- verification must be cheap enough to run on every settled question.

\begin{figure}[htbp]
\centering
\begin{tikzpicture}[
  font=\sffamily\scriptsize,
  st/.style={draw=ink0,fill=fillA,rounded corners=1.5pt,minimum height=9mm,
             minimum width=27mm,align=center,line width=0.5pt},
  dec/.style={draw=ink0,fill=white,shape=diamond,aspect=1.7,inner sep=1pt,
              align=center,line width=0.5pt,font=\sffamily\tiny},
  resx/.style={draw=ink2,fill=fillB,rounded corners=1.5pt,minimum height=8mm,
              minimum width=25mm,align=center,line width=0.4pt},
  ar/.style={-{Latex[length=1.3mm,width=1.1mm]},draw=ink0,line width=0.5pt},
  lbl/.style={font=\sffamily\tiny,color=ink2,inner sep=1.5pt},
]
\node[st] (in) {peer answers\\settled};
\node[dec,right=8mm of in] (cheap) {checklist\\match?};
\node[resx,above right=5mm and 8mm of cheap] (verdict) {supported\\no call spent};
\node[st,below right=5mm and 8mm of cheap] (clf) {1$\times$C bounded\\classify};
\node[dec,right=8mm of clf] (conf) {confidence\\$\ge$ 0.55?};
\node[resx,right=9mm of conf] (ok) {verdict:\\supported};
\node[resx,below=6mm of conf] (poll) {doubtful $\to$\\poll-first};
\node[resx,below=6mm of poll] (esc) {escalate to L iff\\genuinely torn};

\draw[ar] (in) -- (cheap);
\draw[ar] (cheap) -- node[lbl,above,sloped] {yes} (verdict);
\draw[ar] (cheap) -- node[lbl,below,sloped] {no} (clf);
\draw[ar] (clf) -- (conf);
\draw[ar] (conf) -- node[lbl,above,sloped] {yes} (ok);
\draw[ar] (conf) -- node[lbl,right] {no} (poll);
\draw[ar] (poll) -- (esc);
\end{tikzpicture}
\caption{Auditor verification order. The cheap path is attempted first: a checklist match
settles the question with no model call at all. Confidence below the configured threshold
yields a poll rather than a counter-assertion.}
\label{fig:auditor}
\end{figure}

Three properties of this design are worth stating explicitly.

\paragraph{Cheap path first.} Peer answers are compared against the claim checklist of the
retrieved chunks. Where a claim is directly supported or directly contradicted, the verdict
costs nothing. Only disagreement that the checklist cannot settle reaches the model.

\paragraph{Poll before assertion.} When the Auditor is doubtful, it does not override the
peer. It opens a consensus poll on the contested claim, and the group votes. This
preserves the design position --- the circle decides, the AI structures the decision ---
and it converts an AI judgement into a collective one, which is also the pedagogically
better outcome.

\paragraph{Escalation is explicit.} A low compact-tier confidence sets an escalation flag
rather than silently invoking the larger model. The flag is what the router reads, so the
larger-tier budget is spent by a recorded decision with a recorded reason.

\section{The learning-DNA passport}
\label{sec:dna}

Personalization in this system is a small number of well-chosen questions and a long
half-life for their answers. The diagnostic is four items, designed to complete in under
three minutes on a phone, and each item changes a specific downstream behaviour.

\begin{table}[htbp]
\centering
\caption{The four diagnostic items and the state each one writes.}
\label{tab:dna}
\footnotesize
\begin{tabularx}{\textwidth}{@{}p{0.9cm}p{3.6cm}p{2.5cm}X@{}}
\toprule
\textbf{Item} & \textbf{Prompt (trilingual)} & \textbf{Writes} & \textbf{Effect} \\
\midrule
L1 & Level in the current topic: novice / beginner / intermediate
   & \code{level}, \code{levelByTopic} & Chunk level filter; ladder starting point;
     difficulty gate \\
S1 & What helps most: a picture, a daily-life analogy, or numbered steps
   & \code{styleWeights} & Initial weights over six styles; retrieval variant selection \\
G1 & Preferred register: Arabic, Darija, French, English
   & \code{lang} & Response register, chunk language filter, UI locale \\
C1 & Confidence: sure / unsure / lost
   & \code{hintEagerness} & Ladder eagerness ($0.3 / 0.6 / 0.9$); invite flag \\
\bottomrule
\end{tabularx}
\end{table}

The style dimensions are six rather than three (\code{visual}, \code{analogy},
\code{steps}, \code{example}, \code{repetition}, \code{practice}) because the diagnostic
can only \emph{anchor} a prior, not determine a learner. After each check-question the
bandit performs one documented update.

\begin{figure}[htbp]
\centering
\begin{tikzpicture}[
  font=\sffamily\scriptsize,
  ar/.style={-{Latex[length=1.3mm,width=1.1mm]},draw=ink0,line width=0.5pt},
  bx/.style={draw=ink2,fill=white,rounded corners=1.5pt,inner sep=4pt,align=center,
             line width=0.4pt},
  hl/.style={draw=ink0,fill=fillB,rounded corners=1.5pt,inner sep=4pt,align=center,
             line width=0.5pt},
]
\node[bx] (start) {even weights\\each $1/6$};
\node[hl,right=10mm of start] (anchor) {anchor S1:\\$w_S \mathrel{+}= 0.25$\\renormalise};
\node[bx,right=10mm of anchor] (check) {check-question\\outcome $s\in\{0,1\}$};
\node[bx,right=10mm of check] (upd) {$w \mathrel{+}= 0.1\,(s-0.5)$\\renormalise};
\node[bx,right=10mm of upd] (pick) {argmax $w$\\$\Rightarrow$ style};
\draw[ar] (start) -- (anchor);
\draw[ar] (anchor) -- (check);
\draw[ar] (check) -- (upd);
\draw[ar] (upd) -- (pick);
\end{tikzpicture}
\caption{Style adaptation. The passport supplies a prior; a single documented bandit step
per check-question moves it. Weights always sum to one and never fall below a floor, so a
style cannot be permanently extinguished by early bad luck.}
\label{fig:bandit}
\end{figure}

The update is $\alpha = 0.1$ against a baseline of $0.5$, with a floor of $0.01$ per style
followed by renormalisation. The floor matters: without it, a learner who has two unlucky
visual explanations in a row could drive \code{visual} near zero and never see a diagram
again, which is a worse error than the one the bandit was correcting.

\section{The Whisper relay (teach to learn)}
\label{sec:whisper}

Explaining is the highest-value learner activity in the system, and it is also where
misconceptions become visible. The Whisper relay asks a learner to paraphrase a chain of at
least two links, then compares the paraphrase against the claim checklist of the source
chunks and reports a \emph{drift map} with four categories:

\begin{itemize}[leftmargin=1.4em,itemsep=1pt]
  \item \textbf{kept} --- claims reproduced correctly;
  \item \textbf{lost} --- claims present in the source and absent from the paraphrase;
  \item \textbf{distorted} --- claims reproduced with a changed meaning;
  \item \textbf{added} --- claims asserted that the sources do not support.
\end{itemize}

A T0 comparison runs first and settles obvious cases; only a genuinely ambiguous chain
reaches a compact call. The chain must have at least two links, which is validated
server-side rather than trusted from the client --- a one-link ``chain'' is a learner
copying a sentence, not explaining it. After the drift map, the learner gets exactly one
repair round: the pedagogic value is in repairing the distortion, not in being graded on it.

\section{Engineered behaviour under limited capacity}
\label{sec:degradation}

This is the section the design position most directly determines. If the AI must not
replace the circle, then losing the AI must not end the session. The system therefore
models availability as four states, and each state changes what the interface
\emph{promises}.

\begin{figure}[htbp]
\centering
\begin{tikzpicture}[
  font=\sffamily\scriptsize,
  st/.style={draw=ink0,minimum height=13mm,minimum width=27mm,align=center,
             line width=0.5pt,rounded corners=1.5pt},
  g/.style={st,fill=white},
  o/.style={st,fill=fillA},
  r/.style={st,fill=fillB},
  b/.style={st,fill=fillC,text=white},
  ar/.style={-{Latex[length=1.4mm,width=1.2mm]},draw=ink0,line width=0.6pt},
  back/.style={-{Latex[length=1.2mm,width=1.0mm]},draw=ink2,line width=0.4pt,dashed,bend left=24},
  lbl/.style={font=\sffamily\tiny,color=ink2,align=center,inner sep=1.5pt},
]
\node[g] (green) {\textbf{GREEN}\\ normal\\ C and L};
\node[o,right=14mm of green] (orange) {\textbf{ORANGE}\\ L queued / slow\\ C continues};
\node[r,right=14mm of orange] (red) {\textbf{RED}\\ L quota out\\ C only};
\node[b,right=14mm of red] (black) {\textbf{BLACK}\\ provider down\\ offline pack};

\draw[ar] (green) -- node[lbl,above] {L slow} node[lbl,below] {retry $\to$ queue} (orange);
\draw[ar] (orange) -- node[lbl,above] {quota out} node[lbl,below] {C-only labels} (red);
\draw[ar] (red) -- node[lbl,above] {outage} node[lbl,below] {offline pack} (black);
\draw[back] (orange.north) to node[lbl,above] {recovered} (green.north);
\draw[back] (red.south) to node[lbl,below] {quota reset} (orange.south);
\end{tikzpicture}
\caption{The degradation machine. Every transition emits a UI banner, a telemetry event and
a metric gauge, and every state is rehearsable through a development-only chaos toggle.}
\label{fig:degradation}
\end{figure}

\begin{table}[htbp]
\centering
\caption{What each degradation state changes, stated as a product promise.}
\label{tab:degstates}
\footnotesize
\begin{tabularx}{\textwidth}{@{}p{1.55cm}X p{4.2cm}@{}}
\toprule
\textbf{State} & \textbf{Promise kept} & \textbf{Mechanism} \\
\midrule
GREEN  & Full quality, both tiers, citations, capsule polishing
       & Normal routing \\
ORANGE & Peers still first and the session still works; the deeper answer arrives with a
         position and an ETA & L requests enter the FIFO queue with position and ETA
         surfaced; one guided retry, then C fallback \\
RED    & Every question still gets a grounded compact-tier answer with simplified labels;
         no feature disappears silently & L reservations refused atomically; the router
         downgrades; labels and prose are simplified \\
BLACK  & The circle continues with peer work, polls, role-play and flashcards from a
         30-snippet offline pack; nothing is lost & Offline pack served with an
         \code{offline:true} label; writes queued for later sync \\
\bottomrule
\end{tabularx}
\end{table}

\subsection{The quota governor}

The governor is the component that makes the state machine real rather than advisory. It
maintains two ledgers --- per session and per circle per day --- and reserves atomically,
because the obvious implementation is wrong in a way that only appears under load.

\begin{lstlisting}[language=Java,caption={Reserve then compare. Checking before incrementing lets N replicas each read ``under cap'' and all spend.}]
// Increment first, then compare: check-then-increment lets N replicas each read
// "under the cap" and all spend, which is exactly the failure this exists to prevent.
const reservedField = tier === 'C' ? '$reservedC' : '$reservedL';
const capField     = tier === 'C' ? '$cCap'      : '$lCap';
\end{lstlisting}

A single circle can open many sessions in a day, so a per-session cap alone does not bound
a class: ten squads of four against a 20-call session cap is already 200 calls. The
per-circle daily ceiling (\code{CIRCLE\_DAILY\_CALLS}, default 400) is the binding limit,
and it sits above normal classroom use and below a runaway loop.

\subsection{Semantic cache}

Cache keys are a SHA-1 digest over the normalised intent, the level, the language, the
sorted chunk identifiers and the ladder state. The normaliser strips case, diacritics and a
multilingual stop-word list,
so equivalent re-asks collapse to one key and are served without a model call. This is the
cheapest possible win on the ``fewer, well-built requests'' objective, and it is metered:
\code{cacheHitRate} is a first-class metric, not an internal counter.

\subsection{Latency UX}
\label{sec:latency-ux}

Because responses arrive whole rather than streamed, the interface cannot show a partial
answer, and faking a token stream would misrepresent what the system knows. Instead:

\begin{enumerate}[leftmargin=1.6em,itemsep=2pt]
  \item a zero-model-call skeleton renders in under 200\,ms, so the section is never empty;
  \item the compact tier fills the sections it owns;
  \item the larger tier, when it was escalated, slides in afterwards and is labelled
        \emph{deeper answer ready}.
\end{enumerate}

Progress inside a single turn is communicated as the ordered pipeline stages
(\S\ref{sec:cfas}), so the learner sees \emph{retrieving} then \emph{verifying} rather than
a spinner with no semantics.
